%! Function Composition and Decomposition — Unit 1, Lesson 1.3 — Slides (Beamer)
% Order: hook -> I Do -> We Do -> You Do -> debrief -> homework launch (LESSON_SHAPE.md §1).
% \CourseName is NOT defined in beamer — "Precalculus" is written literally.
\documentclass[aspectratio=169, 11pt]{beamer}
\usepackage{precalculus-beamer}   % provides \plumheader{} and \sectionlabel[color]{}

\begin{document}

% TITLE SLIDE
{
\setbeamercolor{background canvas}{bg=plum}
\begin{frame}[plain]
  \vspace{2.8cm}\hspace{0.55cm}%
  \begin{minipage}{0.9\textwidth}
    {\color{goldacc}\bfseries\small LESSON 1.3}\\[0.5cm]
    {\color{white}\bfseries\LARGE Function Composition and Decomposition}\\[1.2cm]
    {\color{lilacmid}\small Precalculus~$\cdot$~Unit 1: Functions and Change}
  \end{minipage}
\end{frame}
}

% HOOK SLIDE
\begin{frame}[t, plain]
  \plumheader{Hook: Two Registers, One Jacket}
  \sectionlabel[goldacc]{Vote Before You Compute}

  A \$60 jacket, a \$10 coupon, 6\% sales tax. Register 1 scans the coupon first,
  then the tax. Register 2 scans the tax first, then the coupon.

  \bigskip
  \begin{columns}[T, onlytextwidth]
    \column{0.46\textwidth}
      \begin{block}{The question}
        Same jacket, same coupon, same tax rate. Do the two of you pay the same
        amount? Commit before anyone computes.
      \end{block}

    \column{0.50\textwidth}
      \begin{block}{Then compute --- and watch the middle number}
        \begin{tabular}{@{}lcc@{}}
          & \textbf{screen} & \textbf{total} \\
          Register 1 & \$50 & \$53.00 \\
          Register 2 & \$63.60 & \$53.60 \\
        \end{tabular}

        \smallskip
        The totals differ because the \textbf{middle numbers} differ.
      \end{block}
  \end{columns}
\end{frame}

% I DO — THE HANDOFF
\begin{frame}[t, plain]
  \plumheader{The Handoff --- The Number In Between}
  \sectionlabel[goldacc]{Notes row THE HANDOFF $\cdot$ I Do: problem 1 $\cdot$ We Do: problem 2}

  \((f \circ g)(x) = f\big(g(x)\big)\), read ``\(f\) of \(g\) of \(x\).'' The circle is
  \textbf{not} multiplication, and the letter written \textbf{last runs first}.

  \medskip
  \begin{center}
  \begin{tikzpicture}[x=1cm, y=1cm, >=stealth]
    \node[draw=plum, very thick, fill=lilac, rounded corners=3pt, minimum width=1.2cm,
          minimum height=0.9cm] (E) at (0,0) {\large \(E\)};
    \node[draw=plum, very thick, fill=lilac, rounded corners=3pt, minimum width=1.2cm,
          minimum height=0.9cm] (T) at (5.2,0) {\large \(T\)};
    \node[font=\small] (in) at (-2.3,0) {\(t\) minutes};
    \node[font=\small] (out) at (7.7,0) {\(T(E(t))\) degrees};
    \draw[->, very thick, plum] (in) -- (E);
    \draw[->, very thick, goldacc] (E) -- (T);
    \draw[->, very thick, plum] (T) -- (out);
    \node[above, font=\small] at (2.6,0.14) {\(E(t)\) feet};
    \node[below, font=\small\bfseries, text=goldacc] at (2.6,-0.10) {THE HANDOFF};
  \end{tikzpicture}
  \end{center}

  \medskip
  \begin{columns}[T, onlytextwidth]
    \column{0.48\textwidth}
      \begin{block}{The cable car}
        \(E(t) = 200t\) feet after \(t\) minutes; \(T(e) = 68 - \frac{e}{100}\) degrees at
        elevation \(e\). \(T(E(4))\): handoff \(\mathbf{800}\) feet, then \(T(800) = 60^\circ\).
      \end{block}

    \column{0.48\textwidth}
      \begin{block}{One number, two jobs}
        The handoff is the inner function's \textbf{answer} and the outer function's
        \textbf{question} at the same time. Find it, and \textbf{write it down}.
      \end{block}
  \end{columns}
\end{frame}

% I DO / WE DO — FROM A TABLE
\begin{frame}[t, plain]
  \plumheader{Finding the Handoff in a Table}
  \sectionlabel[goldacc]{Notes row FINDING IT $\cdot$ I Do: problem 3 $\cdot$ We Do: problems 4--6}

  \begin{columns}[T, onlytextwidth]
    \column{0.44\textwidth}
      \begin{center}
      \renewcommand{\arraystretch}{1.3}
      \begin{tabular}{|l|c|c|c|c|c|}
      \hline
      \(x\) & 0 & 1 & 2 & 3 & 4 \\
      \hline
      \(f(x)\) & 3 & 4 & 0 & 2 & 1 \\
      \hline
      \(g(x)\) & 2 & 3 & 7 & 0 & 4 \\
      \hline
      \end{tabular}
      \end{center}

      \medskip
      \(f(g(1))\): handoff \(g(1) = \mathbf{3}\), then \(f(3) = 2\).

      \smallskip
      \(g(f(1))\): handoff \(f(1) = \mathbf{4}\), then \(g(4) = 4\).

    \column{0.52\textwidth}
      \begin{block}{Swap the machines, swap the handoff}
        Same two functions, same input, a different middle number --- so a different
        answer. Composition is \textbf{not commutative}.
      \end{block}

      \medskip
      \begin{block}{And sometimes there is no answer at all}
        \(f(g(2))\): the handoff is \(7\), and the table has no \(f(7)\). A handoff must be a
        question the \textbf{outer} function accepts --- that is the \textbf{domain of a
        composite}.
      \end{block}
  \end{columns}
\end{frame}

% WE DO — BUILDING THE COMPOSITE
\begin{frame}[t, plain]
  \plumheader{Building the Composite from Formulas}
  \sectionlabel[goldacc]{Notes row BUILDING IT $\cdot$ I Do: problem 7 $\cdot$ We Do: problems 8--10}

  \begin{block}{The recipe}
    Write \(f\) with an empty slot, \(f(\square) = 3(\square) - 2\); put \(g(x)\) in
    \textbf{every} slot; simplify. The handoff is no longer a number --- it is an
    \textbf{expression}. Nothing else changes.
  \end{block}

  \medskip
  \begin{columns}[T, onlytextwidth]
    \column{0.50\textwidth}
      \(f(x) = 3x - 2\), \quad \(g(x) = x^2 + 1\):
      \begin{align*}
        f(g(x)) &= f(x^2+1) = 3(x^2+1) - 2 \\
                &= 3x^2 + 1 \\[4pt]
        g(f(x)) &= g(3x-2) = (3x-2)^2 + 1 \\
                &= 9x^2 - 12x + 5
      \end{align*}

    \column{0.46\textwidth}
      \begin{block}{Two routes, one number}
        \(f(g(2))\) step by step: handoff \(g(2) = 5\), so \(f(5) = 13\).

        \smallskip
        Straight from the rule: \(3(2)^2 + 1 = 13\).
      \end{block}

      \smallskip
      \textbf{Not} \(f(x) \cdot g(x)\): a product has no handoff.
  \end{columns}
\end{frame}

% WE DO — THE CRUX
\begin{frame}[t, plain]
  \plumheader{Two Orders, Two Handoffs --- The Crux}
  \sectionlabel[redacc]{Notes row TWO HANDOFFS $\cdot$ I Do: problem 11 $\cdot$ problem 12 is the crux}

  Coupon \(c(p) = p - 10\), \qquad tax \(t(x) = 1.06x\).

  \medskip
  \begin{columns}[T, onlytextwidth]
    \column{0.48\textwidth}
      \begin{block}{At \(p = 60\)}
        \(t(c(60)) = t(\mathbf{50}) = \$53.00\)

        \smallskip
        \(c(t(60)) = c(\mathbf{63.60}) = \$53.60\)
      \end{block}

      \smallskip
      \begin{block}{At every price}
        \(t(c(p)) = 1.06p - 10.60\) \quad against \quad \(c(t(p)) = 1.06p - 10\)
      \end{block}

    \column{0.48\textwidth}
      \begin{block}{The crux --- problem 12}
        \(f(x) = 2x\), \(g(x) = x - 6\), one input, both orders:

        \smallskip
        \(f(g(10))\): handoff \(\mathbf{4}\), answer \(\mathbf{8}\).

        \smallskip
        \(g(f(10))\): handoff \(\mathbf{20}\), answer \(\mathbf{14}\).
      \end{block}
  \end{columns}

  \bigskip
  \begin{center}
    \textbf{Swap the machines and you swap the handoff. Taxing first charges 6\% of the ten
    dollars you never paid --- at every price.}
  \end{center}
\end{frame}

% WE DO — DECOMPOSITION
\begin{frame}[t, plain]
  \plumheader{Backwards: Decomposition}
  \sectionlabel[redacc]{Notes row TWO HANDOFFS $\cdot$ We Do: problems 13--14}

  \begin{block}{Decomposition = name the handoff}
    Ask what a calculator would compute \textbf{first}. That value is the handoff, and its rule
    is the \textbf{inner} function; whatever is done to it is the \textbf{outer} function.
  \end{block}

  \bigskip
  \begin{columns}[T, onlytextwidth]
    \column{0.48\textwidth}
      \(h(x) = \sqrt{x + 5}\)

      \smallskip
      You would add 5 first, so:

      \smallskip
      Inner \(g(x) = x + 5\), \quad outer \(f(u) = \sqrt{u}\).

    \column{0.48\textwidth}
      \begin{block}{Always check}
        Substitute back: \(f(g(x)) = \sqrt{x+5} = h(x)\). \checkmark

        \smallskip
        A split that does not rebuild \(h\) exactly is the wrong split.
      \end{block}
  \end{columns}
\end{frame}

% YOU DO
\begin{frame}[t, plain]
  \plumheader{You Do --- Table \& Formulas}
  \sectionlabel[redacc]{Notes row TABLE \& FORMULAS $\cdot$ problems 15--18, alone, 14 minutes}

  \begin{columns}[T, onlytextwidth]
    \column{0.40\textwidth}
      \begin{center}
      \renewcommand{\arraystretch}{1.3}
      \begin{tabular}{|l|c|c|c|c|c|}
      \hline
      \(x\) & 1 & 2 & 3 & 4 & 5 \\
      \hline
      \(p(x)\) & 5 & 3 & 1 & 4 & 2 \\
      \hline
      \(q(x)\) & 4 & 1 & 2 & 5 & 3 \\
      \hline
      \end{tabular}
      \end{center}

      \medskip
      \begin{center}
        \(f(x) = x + 3\) \qquad \(g(x) = x^2\)
      \end{center}

    \column{0.56\textwidth}
      \begin{enumerate}\setlength{\itemsep}{7pt}
        \setcounter{enumi}{14}
        \item From the table: \(p(q(2))\) and \(q(p(2))\). Write each handoff. Equal?
        \item Build \((f \circ g)(x)\) and \((g \circ f)(x)\), simplified. Same function?
        \item Find \(f(g(-2))\) by naming the handoff, then check it against your rule from 16.
        \item Decompose \(m(x) = (2x+1)^3\). Which rule runs first when you compute \(m(1)\)?
              Find \(m(1)\).
      \end{enumerate}
  \end{columns}
\end{frame}

% DEBRIEF
\begin{frame}[t, plain]
  \plumheader{Debrief --- The Headline}
  \sectionlabel[goldacc]{Share out 15--18, one answer each $\cdot$ then say it back}

  \begin{columns}[T, onlytextwidth]
    \column{0.48\textwidth}
      \begin{block}{The handoff}
        One number sits between the two machines: the inner function's \textbf{answer} and the
        outer function's \textbf{question}. Find it first and write it down.
      \end{block}
    \column{0.48\textwidth}
      \begin{block}{The order}
        \textbf{Swap the machines and you swap the handoff.} That is why \(f(g(x))\) and
        \(g(f(x))\) are usually different functions.
      \end{block}
  \end{columns}

  \bigskip
  \textbf{Next time (1.4, Inverse Functions):} we hunt for the partner that \textit{undoes} a
  function --- and the test is a composition run in \textbf{both} directions.

  \smallskip
  \textit{And later, in 1.5, a horizontal shift will turn out to be a composition too.}
\end{frame}

% HOMEWORK LAUNCH
\begin{frame}[t, plain]
  \plumheader{Start Your Homework}
  \sectionlabel[redacc]{Homework Launch $\cdot$ 3 minutes, alone, silent}

  \begin{itemize}\setlength{\itemsep}{8pt}
    \item Copy the due date into the \textbf{Homework} row of your cover sheet:
          \quad Due: \rule{3cm}{0.4pt}
    \item Start \textbf{problem 1} and \textbf{problem 2} now. Write the handoff on every line
          before you use it.
    \item The \textbf{Homework} (last two pages) is graded.
    \item \textbf{AP Practice} (the two pages before it) is \textbf{extra credit} --- optional,
          AP-exam style. Do the homework first.
  \end{itemize}
\end{frame}

\end{document}
